\documentclass[11pt,a4paper]{article}
\usepackage[utf8]{inputenc}
\usepackage[T1]{fontenc}
\usepackage{amsmath,amssymb}
\usepackage{booktabs}
\usepackage{hyperref}
\usepackage{geometry}
\geometry{margin=2.5cm}
\usepackage{enumitem}
\usepackage{graphicx}

\title{\textbf{1\_parameters\_estimation} \\ Alonso's Transmittance Model Fitting}
\author{AzONetV2}
\date{}

\begin{document}
\maketitle

\section{Overview}

This folder contains the fitting of experimental transmittance spectra to the theoretical model proposed by Alonso et al. (IIM --- Interference Induced Method) using three global optimization algorithms from SciPy. The production version is in {\tt Python\_Notebooks/}, using the \textbf{AmaroX} custom library.

\section{Context}

In the previous folder ({\tt 0\_data\_processing}), a DataFrame was created containing experimental spectra and their thicknesses. The goal here is to fit each experimental transmittance curve with the theoretical model by minimizing the Mean Square Error (MSE).

\section{Transmittance Model}

The Alonso model describes transmittance as a function of \textbf{13 parameters}:

\begin{table}[htbp]
\centering
\begin{tabular}{ll}
\toprule
\textbf{Parameter} & \textbf{Description} \\
\midrule
$x$            & Wavelength [nm] \\
$d$            & Film thickness [nm] \\
$R_1$, $R_2$   & Surface roughness parameters \\
$A$, $B$, $C$, $D$, $E$ & Sellmeier coefficients (background refractive index) \\
$\alpha$, $\beta$, $\gamma$ & Absorption coefficients \\
$n_e$          & Free carrier concentration \\
\bottomrule
\end{tabular}
\end{table}

The model incorporates:
\begin{itemize}
    \item \textbf{Sellmeier equation} for the background dielectric constant:
    \begin{equation*}
        \varepsilon_{1b}(x) = A + \frac{B x^2}{x^2 - C^2} + \frac{D x^2}{x^2 - E^2}
    \end{equation*}
    \item \textbf{Drude model} for free-carrier contribution (depends on $n_e$):
    \begin{equation*}
        \varepsilon_{1f} = -\frac{3182.61 \cdot n_e}{\omega^2 + \gamma^2}, \quad
        \varepsilon_{2f} = \frac{3182.61 \cdot n_e \cdot \gamma}{\omega(\omega^2 + \gamma^2)}
    \end{equation*}
    where $\omega = 2\pi c \cdot 10^9 / x$ and $\gamma = 2.8 \times 10^{11} \cdot x$.
    \item \textbf{Absorption model} with Urbach tail:
    \begin{equation*}
        \alpha_{\text{total}} = \kappa \frac{4\pi}{x} + \alpha_0 \exp\!\left[1240 \cdot \beta \left(\frac{1}{x} - \frac{1}{\lambda_g}\right)\right]
    \end{equation*}
    \item \textbf{Roughness corrections} for interfaces via scalar scattering theory.
    \item \textbf{Glass substrate} transmittance ({\tt TexpglassO.txt}) to compute $n_g$.
\end{itemize}

The refractive index is constrained to the \textbf{physically meaningful range}: $1.8 \leq n(\lambda) \leq 2.1$.

Physical constants used: $c = 3 \times 10^8$ m/s (speed of light), $\mu = 3.90 \times 10^{-4}$ m$^2$/Vs (mobility).

\section{Data}

\begin{itemize}
    \item \textbf{Source}: {\tt ../../results/dataframe\_spectrum\_thickness\_145\_final.pkl}
    \item \textbf{Samples}: 145 experimental thin-film samples (AZO --- Aluminium-doped Zinc Oxide)
    \item \textbf{Spectrum range}: 190--1100 nm (911 points per spectrum)
    \item \textbf{Glass transmittance}: {\tt ../../experimental\_samples/Background\_data/TexpglassO.txt}
\end{itemize}

\section{Optimization Algorithms}

Three global optimization algorithms were tested, all minimizing the Mean Square Error (MSE):
\begin{equation*}
    \text{MSE} = \frac{1}{N}\sum_{i=1}^{N} \big(T_{\text{model}}(\lambda_i; \Theta) - T_{\text{exp}}(\lambda_i)\big)^2
\end{equation*}

\subsection{1. Basin-Hopping ({\tt basin\_hopping})}
\begin{itemize}
    \item \textbf{File}: {\tt 1\_SciPy\_AllSamples\_BH\_NonLinear\_Porcentual\_145F.ipynb}
    \item \textbf{Local optimizer}: SLSQP
    \item \textbf{Constraint}: {\tt NonlinearConstraint} on refractive index ($1.8 \leq n \leq 2.1$)
    \item \textbf{Initial guess}: Mean values from prior EGP-GA results:
    \begin{itemize}
        \item $R_1 = 30.6$, $R_2 = 26.5$
        \item $A = 2.0065$, $B = 1.574 \times 10^6$, $C = 10^7$, $D = 1.5868$, $E = 260.63$
        \item $\alpha_0 = 5.0 \times 10^{-3}$, $\beta = 10$, $\lambda_g = 360$
        \item $n_e = 2.1 \times 10^{26}$
    \end{itemize}
    \item \textbf{Bounds}: Allowed $\pm 50\%$ variation around seed values ($pl = 0.5$), $\pm$ measurement error for thickness.
\end{itemize}

\subsection{2. Differential Evolution ({\tt differential\_evolution})}
\begin{itemize}
    \item \textbf{File}: {\tt 2\_SciPy\_AllSamples\_DE\_NonLinear\_Porcentual\_145F.ipynb}
    \item \textbf{Constraint}: {\tt NonlinearConstraint} on refractive index ($1.8 \leq n \leq 2.1$)
    \item \textbf{Parallelization}: {\tt workers=-1} (uses all available CPUs)
    \item \textbf{Stochastic}: Random seed per run ({\tt random.randint(0, 2**32 - 1)})
    \item \textbf{Bounds}: Same $\pm 50\%$ around seeds, $\pm$ measurement error for thickness.
\end{itemize}

\subsection{3. Direct ({\tt direct})}
\begin{itemize}
    \item \textbf{File}: {\tt 3\_SciPy\_AllSamples\_Direct\_NonLinear\_Porcentual\_145F.ipynb}
    \item \textbf{Constraint}: Penalty method --- Direct does not natively support {\tt NonlinearConstraint}. A wrapper function {\tt new\_error} adds a penalty term when the refractive index constraint is violated:
    \begin{verbatim}
    if (refractive_index(params) > 2.1) or (refractive_index(params) < 1.8):
        _error += _error + refractive_index(params)
    \end{verbatim}
    \item \textbf{Fastest} but weakest results.
\end{itemize}

\section{Results}

Results are saved as NumPy arrays in {\tt ../../results/SciPy\_IIM/}:

\begin{table}[htbp]
\centering
\begin{tabular}{ll}
\toprule
\textbf{File} & \textbf{Algorithm} \\
\midrule
{\tt basin\_hopping\_NonLinear\_145F.npy} & Basin-Hopping \\
{\tt differential\_evolution\_NonLinear\_145F.npy} & Differential Evolution \\
{\tt direct\_NonLinear\_145F.npy} & Direct \\
\bottomrule
\end{tabular}
\end{table}

Each array has shape $(145, 14)$ with columns:
\begin{verbatim}
[SampleIndex, Thickness, R1, R2, A, B, C, D, E, alpha, beta, gamma, ne, MSE]
\end{verbatim}

\section{Comparison \& Conclusions}

The notebook {\tt 4\_Comparison.ipynb} (which also generates {\tt comparativa\_optimizadores.pdf}) compares the three optimizers:

\begin{table}[htbp]
\centering
\begin{tabular}{lccc}
\toprule
\textbf{Optimizer} & \textbf{Mean MSE} & \textbf{Execution Time} & \textbf{Quality} \\
\midrule
\textbf{Differential Evolution} & $\sim 0$ (best) & $\sim 4$ hours & Excellent \\
Basin-Hopping & $\sim 5$ & $\sim 12$ hours & Moderate \\
Direct & $\sim 8$--$30$ & $\sim 5$ minutes & Poor \\
\bottomrule
\end{tabular}
\end{table}

\textbf{Key findings:}
\begin{itemize}
    \item \textbf{Differential Evolution} is the best optimizer for this problem --- it achieves the lowest MSE consistently across all samples. Its KDE density distribution is tightly centered near zero.
    \item \textbf{Basin-Hopping} produces a normal-like distribution centered around MSE $\approx 5$, with significantly longer runtime.
    \item \textbf{Direct} is fast but produces large errors, with a multi-modal error distribution. Not recommended for production use.
\end{itemize}

The comparison visualizations are saved in the {\tt images/} folder:
\begin{itemize}
    \item {\tt comparative.png} --- Scatter plot of MSE per sample for all three algorithms
    \item {\tt kde\_optimizers.png} --- Kernel Density Estimation of the MSE distributions
    \item {\tt final\_comparative.png} --- Bar chart of mean MSE comparison
\end{itemize}

The PDF report {\tt comparativa\_optimizadores.pdf} (generated via {\tt fpdf}) includes Spanish-language commentary alongside these figures.

\section{Python\_Notebooks Subfolder}

Contains {\tt .py} script versions of the three optimization notebooks, refactored to use the \textbf{AmaroX} custom library instead of inline model definitions:

\begin{table}[htbp]
\centering
\begin{tabular}{ll}
\toprule
\textbf{Script} & \textbf{Equivalent Notebook} \\
\midrule
{\tt 1\_SciPy\_AllSamples\_BH\_NonLinear\_Porcentual\_145F.py} & Notebook 1 (Basin-Hopping) \\
{\tt 2\_SciPy\_AllSamples\_DE\_NonLinear\_Porcentual\_145F.py} & Notebook 2 (Differential Evolution) \\
{\tt 3\_SciPy\_AllSamples\_Direct\_NonLinear\_Porcentual\_145F.py} & Notebook 3 (Direct) \\
\bottomrule
\end{tabular}
\end{table}

These scripts import the {\tt modelo\_transmitancia} function and other utilities from {\tt AmaroX}, making them cleaner for production use. The model equations, constraints, bounds, and output paths are identical to their notebook counterparts.

\section{Test-400nm-1100nm Subfolder}

This subfolder explores the effect of \textbf{restricting the fitting range} from the full spectrum (190--1100 nm) to the visible--NIR region (400--1100 nm). The low-wavelength UV tail (190--400 nm) is noisy and dominated by absorption.

\subsection{Adjustment Notebook}
\textbf{File:} {\tt Test-400nm-1100nm/Adjustment\_400nm\_1100nm.ipynb}

Reruns \textbf{Differential Evolution} fitting on all 145 samples, but the MSE error function only considers wavelengths $\geq 400$ nm ({\tt y\_pred[210:]} in the 911-point array). The model, constraints, bounds, and glass data are identical to the main DE notebook. Results are saved as {\tt differential\_evolution\_NonLinear\_145F.npy} in the subfolder.

\subsection{Comparison Notebook}
\textbf{File:} {\tt Test-400nm-1100nm/Compare\_Old\_New\_Fittings.ipynb}

Compares the two fitting strategies across all 145 spectra:

\begin{table}[htbp]
\centering
\begin{tabular}{lcc}
\toprule
\textbf{Metric} & \textbf{Old (190--1100 nm)} & \textbf{New (400--1100 nm)} \\
\midrule
Mean MSE      & $1.099$    & $1.186$ \\
Median MSE    & $0.757$    & $0.743$ \\
Wins (lower MSE) & \textbf{96 / 145 (66.2\%)} & 49 / 145 (33.8\%) \\
\bottomrule
\end{tabular}
\end{table}

\textbf{Takeaway:} Fitting the full 190--1100 nm range is recommended for production use. The 400--1100 nm restricted fit can serve as a sensitivity analysis but does not improve overall fitting quality.

\section{Dependencies}
\begin{itemize}
    \item {\tt numpy}, {\tt pandas}, {\tt matplotlib}, {\tt seaborn}, {\tt scipy}
    \item {\tt fpdf} (for PDF report generation in {\tt 4\_Comparison.ipynb})
    \item {\tt AmaroX} (custom library, used only in {\tt Python\_Notebooks/*.py})
\end{itemize}

\end{document}
